% Generated draft: requires mathematical and source review.
\documentclass[11pt]{article}
\usepackage[margin=25mm]{geometry}
\usepackage{fontspec}
\setmainfont{Latin Modern Roman}
\usepackage{amsmath,amssymb,mathtools,unicode-math}
\setmathfont{Latin Modern Math}
\usepackage{enumitem,needspace,xurl,hyperref}
\setlength{\parindent}{0pt}
\setlength{\parskip}{5pt}
\setlength{\emergencystretch}{3em}
\newcommand{\N}{\mathbb{N}}
\newcommand{\Z}{\mathbb{Z}}
\newcommand{\Q}{\mathbb{Q}}
\newcommand{\R}{\mathbb{R}}
\newcommand{\C}{\mathbb{C}}
\newcommand{\F}{\mathbb{F}}
\newcommand{\A}{\mathbb{A}}
\newcommand{\PP}{\mathbb{P}}
\newcommand{\E}{\mathbb{E}}
\newcommand{\eps}{\varepsilon}
\newcommand{\dd}{\,\mathrm{d}}
\begin{document}
\section*{UnsolvedMath ID: 20000806}
\textbf{Source problem code:} catalogue entry
\subsection{Definitions and mathematical statement}
For a field $k$, $\operatorname{Hilb}^d(\A_k^n)$ parametrizes zero-dimensional subschemes of length $d$. ``A nonreduced component'' is read as generic nonreducedness of the Hilbert scheme along an irreducible component, not nonreducedness of its support given the reduced scheme structure.

\textbf{Problem 5.} Do there exist $d,n$ for which $\operatorname{Hilb}^d(\A^n)$ has a nonreduced component? If so, exhibit one. The source does not fix the field or characteristic.

\textbf{Supplied characteristic-zero answer.} Over an algebraically closed field of characteristic zero, the aid attributes to Jelisiejew's 2024 theorem the following examples. For each $s\in\{6,7,8,9\}$, the very-compressed locus of local algebras with Hilbert function
\[
 (\dim A/\mathfrak m,\dim\mathfrak m/\mathfrak m^2,
 \dim\mathfrak m^2/\mathfrak m^3,\dim\mathfrak m^3/\mathfrak m^4)
 =(1,4,10,s),\qquad\mathfrak m^4=0,
\]
has closure equal to the reduced support of a generically nonreduced irreducible component of $\operatorname{Hilb}^{15+s}(\A^4)$. Here the length is $1+4+10+s=15+s$. The exact theorem's definition of the very-compressed locus and its dense cubic-generated open part must be checked against the underlying reference.

On the dense cubic-generated parts the report gives the lower bounds
\[
 \begin{array}{c|rrrr}
 d&21&22&23&24\\\hline
 \dim T_{[Z]}\operatorname{Hilb}^d(\A^4)-\dim C_{\mathrm{red}}
 &\geq10&\geq7&\geq12&\geq25
 \end{array}
\]
for the corresponding component $C$, at the general points specified by the report. A tangent-excess test applies on a dense smooth locus of the reduced support away from other components: excess there detects generic nonreducedness. Excess at a single arbitrary singular or intersection point would not suffice.

\textbf{Reported propagation.} If $Z$ is one of the finite subschemes and $p_1,\ldots,p_m$ are distinct reduced points outside its support, then the aid proves the completed-local-ring statement
\[
 \widehat{\mathcal O}_{\operatorname{Hilb}^{d+m}(\A^4),[Z\sqcup\{p_1,\ldots,p_m\}]}
 \simeq
 \widehat{\mathcal O}_{\operatorname{Hilb}^{d}(\A^4),[Z]}
 [[t_1,\ldots,t_{4m}]].
\]
This noncanonical smooth factor adds $4m$ to both tangent and reduced-component dimensions. Hence the stated lower-bound excess and generic nilpotents persist on the associated component. Starting with length $21$ supplies a generically nonreduced component of $\operatorname{Hilb}^d(\A^4)$ for every $d\geq21$.

The aid therefore reports a solved existence problem in characteristic zero. It does not claim minimality of length $21$, classify all such components, or give the same statement in positive characteristic.
\subsection{Short English statement}
The source asks for a genuinely generically nonreduced component of a Hilbert scheme of points. The aid cites characteristic-zero examples in affine four-space of lengths $21$ through $24$, with explicit tangent-excess lower bounds. Adding disjoint reduced points preserves the completed local structure up to smooth variables and yields every larger length.
\subsection{Sources}
\begin{itemize}
\item Local UnsolvedMath registry, permanent ID $20000806$, \texttt{catalogue entry}; canonical \texttt{aim-arithmetic-geometry-notes.json}, \texttt{notes[51]}.
\item 2010 AIM workshop ``Components of Hilbert schemes,'' Problem 5: \url{https://aimath.org/WWN/hilbertschemes/hilbertschemes.pdf}.
\item Jelisiejew's 2024 theorem, as named in the supplied summary. No title, identifier, theorem number or complete bibliographic reference is supplied.
\item Research attempt 1, \texttt{full\_solution}, run status \texttt{solved}, internal rigor \texttt{proved}, low novelty confidence; registry \texttt{SOLVED-IN-LITERATURE}. Explicitly a machine-generated research aid, not peer reviewed; independently verify before citation. Original literature field empty.
\end{itemize}
\end{document}
